\documentclass[11pt]{article}
\usepackage[T1]{fontenc}
\usepackage{lmodern,textcomp,amsmath,amssymb}
\usepackage[margin=1in]{geometry}
\usepackage[hidelinks]{hyperref}
\setlength{\parindent}{0pt}
\setlength{\parskip}{0.7em}
\pagestyle{empty}
\begin{document}
{\Large\bfseries Manuscript submission}\par
Finite Fields and Their Applications\hfill 26 October 2026

Dear Editors,

Please consider the manuscript \emph{The Welch function is not third-order sum-free} by Deep Bhattacharjee for publication in \emph{Finite Fields and Their Applications}.

The manuscript proves that the Welch function $X^{2^m+3}$ on $\mathbb{F}_{2^{2m+1}}$ is not third-order sum-free for every $m\ge3$, settling the conjecture formulated by Hou and Zhao. It gives a two-trace criterion and the exact count $(2^n+1-3K_n)/4$, with $K_n$ a binary Kloosterman sum. The corresponding number of three-dimensional linear zero-sum subspaces is $N_n(2^n-1)/42$. The proof is written out from the reduction to the final character-sum count and uses the classical Weil bound. The supplied code checks the formulas independently, and the supplied Lean files check specified identities and small cases. The finite-field structure, character sums and relation to APN functions fit the journal\textquotesingle s scope.

An earlier version is publicly available as Hexagon \href{https://hexagonmath.org/2610.00472v1}{2610.00472v1}. The source, verification code and supporting material are archived with the accompanying release:
\begin{center}\url{https://github.com/creelie/eConjectureWelchSumFree/releases/tag/v1.2.0}\end{center}
The submission includes a separate declaration describing AI assistance and the scope of the formal verification.

Thank you for considering the manuscript.

Sincerely,\par
Deep Bhattacharjee\par
Corresponding author\par
\href{mailto:itsdeep@live.com}{itsdeep@live.com}\par
\href{mailto:d.bhattacharjee@erl-forschung.de}{d.bhattacharjee@erl-forschung.de}\par
ORCID: \href{https://orcid.org/0000-0003-0466-750X}{0000-0003-0466-750X}
\end{document}
